\documentclass[11pt,a4paper]{article}
\usepackage[margin=20mm,top=19mm,bottom=20mm]{geometry}
\usepackage{fontspec,amsmath,amssymb,xcolor,enumitem,fancyhdr}
\setmainfont{NotoSansKR-Regular.ttf}[Path=../../../common/assets/fonts/,BoldFont=NotoSansKR-Semibold.ttf]
\setmonofont{Latin Modern Mono}
\definecolor{ink}{HTML}{203B50}
\definecolor{rulegray}{HTML}{D2D9DF}
\pagestyle{fancy}\fancyhf{}
\fancyhead[L]{\small Machine Learning Concept Check}\fancyhead[R]{\small W1}
\fancyfoot[C]{\small\thepage\ / 4}
\renewcommand{\headrulewidth}{0.3pt}
\setlength{\headheight}{15pt}
\setlength{\parindent}{0pt}\setlength{\parskip}{5pt}
\linespread{1.2}
\newcommand{\parttitle}[2]{{\large\bfseries\color{ink} #1}\par{\small #2}\par\vspace{5mm}}
\newcommand{\question}[2]{\textbf{#1.} #2\par}
\newcommand{\answerlines}[1]{\foreachline{#1}}
\newcount\linecount
\newcommand{\foreachline}[1]{\linecount=0\loop\ifnum\linecount<#1\vspace{8mm}\par{\color{rulegray}\hrule}\advance\linecount by1\repeat}
\begin{document}

{\LARGE\bfseries\color{ink} W1 · Machine Learning and Linear Regression}\par
{\small Supervised learning · Regression · Cost functions · Gradient descent}\par
\vspace{3mm}Name\enspace\rule{43mm}{0.3pt}\hfill Date\enspace\rule{35mm}{0.3pt}\par\vspace{3mm}
\parttitle{1. Concept Check}{Choose the single best answer for each question.}
\question{1}{A model sees many pairs of house sizes and sale prices, then predicts the price of a new house. What kind of learning is this?}
\begin{enumerate}[label=(\arabic*),leftmargin=9mm,itemsep=0pt,topsep=2pt,parsep=0pt]
\item Supervised learning
\item Grouping houses without target labels
\item A calculation that uses no data
\item Choosing prices at random
\end{enumerate}\vfill
\question{2}{Which is a regression problem?}
\begin{enumerate}[label=(\arabic*),leftmargin=9mm,itemsep=0pt,topsep=2pt,parsep=0pt]
\item Classifying email as spam or legitimate
\item Classifying animals in photos as cats or dogs
\item Predicting tomorrow's temperature as a number
\item Classifying reviews as positive or negative
\end{enumerate}\vfill
\question{3}{You have customer data without labels and want to group similar customers. Which approach is most appropriate?}
\begin{enumerate}[label=(\arabic*),leftmargin=9mm,itemsep=0pt,topsep=2pt,parsep=0pt]
\item Assigning the same target label to every customer
\item Finding groups with unsupervised learning
\item Using customer IDs directly as purchase amounts
\item Fixing the groups before looking at the data
\end{enumerate}\vfill
\question{4}{In house-price prediction, $x$ is house size and $y$ is the actual price. Which is the input feature?}
\begin{enumerate}[label=(\arabic*),leftmargin=9mm,itemsep=0pt,topsep=2pt,parsep=0pt]
\item The actual price
\item The amount of prediction error
\item The number of training iterations
\item House size
\end{enumerate}\vfill
\question{5}{In $\hat y=wx+b$, you keep $b$ fixed and change only $w$. Which property of the line changes directly?}
\begin{enumerate}[label=(\arabic*),leftmargin=9mm,itemsep=0pt,topsep=2pt,parsep=0pt]
\item Its slope
\item The number of examples
\item The actual targets
\item The input units
\end{enumerate}\vfill
\newpage
\parttitle{1. Concept Check / Continued}{Choose the single best answer for each question.}
\question{6}{If $\hat y=2x+1$ and $x=3$, what does the model predict?}
\begin{enumerate}[label=(\arabic*),leftmargin=9mm,itemsep=0pt,topsep=2pt,parsep=0pt]
\item $3$
\item $5$
\item $7$
\item $9$
\end{enumerate}\vfill
\question{7}{Which best describes the purpose of a cost function?}
\begin{enumerate}[label=(\arabic*),leftmargin=9mm,itemsep=0pt,topsep=2pt,parsep=0pt]
\item Naming the training data
\item Summarizing the differences between predictions and actual values in one number
\item Increasing the number of input features
\item Always producing the correct answer
\end{enumerate}\vfill
\question{8}{You compare two models on the same data using the same cost function. Model A has a lower cost than B. What does this mean?}
\begin{enumerate}[label=(\arabic*),leftmargin=9mm,itemsep=0pt,topsep=2pt,parsep=0pt]
\item A must be more accurate on every new example
\item A has not been trained at all
\item B must have more data
\item A has a smaller summary of prediction errors on this data
\end{enumerate}\vfill
\question{9}{Why does gradient descent use the gradient?}
\begin{enumerate}[label=(\arabic*),leftmargin=9mm,itemsep=0pt,topsep=2pt,parsep=0pt]
\item To adjust parameters step by step in a direction that reduces cost
\item To change actual targets to match predictions
\item To delete training data regardless of its usefulness
\item To keep making the gradient larger
\end{enumerate}\vfill
\question{10}{What can happen if the learning rate is too large?}
\begin{enumerate}[label=(\arabic*),leftmargin=9mm,itemsep=0pt,topsep=2pt,parsep=0pt]
\item The model must reach the minimum in one step
\item The number of input features automatically increases
\item Updates can overshoot the minimum, causing the cost to increase or oscillate
\item Every prediction becomes equal to its actual target
\end{enumerate}\vfill
\newpage
\parttitle{2. Think Like an ML Scientist}{Question 11 · Read the pseudocode and find why the model does not work as intended.}
You built a model that predicts house prices from size. Small parameter updates should reduce the cost, but it increases as training continues. Assume the learning rate is sufficiently small and the gradient function is implemented correctly.
\par\vspace{3mm}
{\small\ttfamily\begin{tabular}{r l}
1 & \hspace*{0.0mm}w = 0; b = 0 \\
2 & \hspace*{0.0mm}alpha = 0.001 \\
3 & \hspace*{0.0mm}repeat 1000 times: \\
4 & \hspace*{6.0mm}dw, db = gradient(X, y, w, b) \\
5 & \hspace*{6.0mm}w = w + alpha * dw \\
6 & \hspace*{6.0mm}b = b + alpha * db \\
\end{tabular}}\par\vspace{5mm}
\textbf{11(1)} Identify the two lines that work against the goal of reducing cost.
\answerlines{5}
\vspace{6mm}
\textbf{11(2)} Explain why the cost can increase and rewrite each incorrect line as one corrected line.
\answerlines{5}
\vspace{6mm}
\newpage
\parttitle{3. Open Question}{Question 12 · There is no single correct answer. Describe your method and reasoning.}
Suppose a tool converts sentences into lists of numbers of equal length. You use it on 100 customer reviews. How would you discover patterns among the reviews?
\par\vspace{3mm}
Describe the pattern you want to find, what you would compare, and how you would check whether it is meaningful. You do not have to assume that numerical distance necessarily measures a difference in meaning.
\par\vspace{3mm}
{\small Equations and code are optional. You may answer with text, a diagram, or a simple flowchart.}
\answerlines{15}
\vfill{\footnotesize These original practice questions are based on the corresponding week of Machine Learning Specialization, Course 1.}
\end{document}
